\section{Một số hướng nghiên cứu gần đây về phương pháp
\protect\tr{kernel}}

Trong những năm gần đây, các phương pháp \tr{kernel} tiếp tục được nghiên cứu
mạnh mẽ theo nhiều hướng khác nhau. Một mặt, các công trình mới tìm cách mở
rộng khả năng tính toán của phương pháp \tr{kernel} cho dữ liệu lớn. Mặt khác,
nhiều nghiên cứu khai thác mối liên hệ giữa phương pháp \tr{kernel} và mạng
nơ-ron sâu, đặc biệt thông qua các khái niệm như \tr{Kernel} Tiếp tuyến Nơ-ron
(Neural Tangent Kernel), Đặc trưng Ngẫu nhiên (Random Features), RKHS và các
mở rộng của RKHS.

Trong phần này, ta khảo sát ba bài báo tiêu biểu sau năm 2021 có liên quan
trực tiếp đến chủ đề phương pháp \tr{kernel}:
\begin{enumerate}
  \item Toward Large Kernel Models;
  \item Random Gegenbauer Features for Scalable Kernel Methods;
  \item Deep Learning with Kernels through RKHM and the Perron--Frobenius
  Operator.
\end{enumerate}

\subsubsection{Toward Large Kernel Models}

Bài báo \emph{Toward Large Kernel Models} được công bố tại ICML 2023. Bài toán
chính mà bài báo đặt ra là làm thế nào để mở rộng phương pháp \tr{kernel} cho
các tập dữ liệu lớn trong khi vẫn giữ được khả năng biểu diễn mạnh của mô hình
\tr{kernel}.

Trong máy \tr{kernel} (kernel machine) truyền thống, hàm dự đoán thường có
dạng:
\begin{equation}
  f(x) = \sum_{i=1}^{n} \alpha_i K(x,x_i),
\end{equation}
phụ thuộc trực tiếp vào số mẫu huấn luyện $n$. Vì vậy, khi dữ liệu lớn, kích
thước ma trận \tr{kernel} và chi phí tính toán tăng rất nhanh. Điều này làm
cho phương pháp \tr{kernel} khó cạnh tranh với học sâu (deep learning) trong
các bài toán quy mô lớn.

Đóng góp chính của bài báo là đề xuất mô hình \tr{kernel} tổng quát
(general kernel model), trong đó số lượng tâm (center) có thể được chọn độc
lập với số lượng mẫu huấn luyện. Mô hình có dạng:
\begin{equation}
  f(x) = \sum_{j=1}^{p} \alpha_j K(x,z_j),
\end{equation}
với $p$ là kích thước mô hình và $\{z_j\}_{j=1}^{p}$ là tập các tâm. Khi đó,
$p$ không nhất thiết phải bằng $n$. Nhờ vậy, kích thước mô hình và kích thước
dữ liệu có thể được mở rộng một cách độc lập.

Bên cạnh đó, bài báo giới thiệu thuật toán EigenPro 3.0, dựa trên hạ gradient
ngẫu nhiên đối ngẫu có tiền điều kiện hóa và phép chiếu (projected dual
preconditioned stochastic gradient descent). Thuật toán này giúp huấn luyện
các mô hình \tr{kernel} lớn hơn nhiều so với các phương pháp \tr{kernel}
truyền thống, đồng thời có thể tận dụng tính toán trên GPU.

Mối liên hệ với lý thuyết \tr{kernel} trong chương đã chọn là rất trực tiếp.
Bài báo sử dụng biểu diễn \tr{kernel} trong RKHS, định lý biểu diễn và bài
toán học hàm trong không gian Hilbert sinh bởi \tr{kernel}. Nếu trong lý
thuyết cơ bản, nghiệm của bài toán \tr{kernel} thường được biểu diễn thông qua
toàn bộ tập huấn luyện, thì bài báo này mở rộng cách biểu diễn đó bằng cách
tách tập tâm \tr{kernel} khỏi tập dữ liệu. Do đó, bài báo có thể được xem là
một hướng phát triển nhằm làm cho lý thuyết \tr{kernel} cổ điển phù hợp hơn
với bài toán học máy quy mô lớn.

\subsubsection{Random Gegenbauer Features for Scalable Kernel Methods}

Bài báo \emph{Random Gegenbauer Features for Scalable Kernel Methods} được
công bố tại ICML 2022. Bài toán chính của bài báo là xây dựng một phương pháp
xấp xỉ \tr{kernel} hiệu quả cho một lớp \tr{kernel} rộng hơn so với các phương
pháp Đặc trưng Ngẫu nhiên truyền thống.

Các phương pháp \tr{kernel} thường gặp khó khăn vì cần lưu trữ và xử lý ma
trận Gram:
\begin{equation}
  K(X,X) = \bigl(K(x_i,x_j)\bigr)_{i,j=1}^{n},
\end{equation}
có kích thước $n \times n$. Khi $n$ lớn, chi phí bộ nhớ và tính toán trở thành
rào cản chính. Đặc trưng Fourier ngẫu nhiên giải quyết vấn đề này cho các
\tr{kernel} bất biến tịnh tiến bằng cách xấp xỉ \tr{kernel} thông qua một ánh
xạ đặc trưng ngẫu nhiên hữu hạn chiều.

Đóng góp chính của bài báo gồm các điểm sau. Thứ nhất, tác giả xây dựng lớp
\tr{Kernel} Generalized Zonal (Generalized Zonal Kernels) và đưa ra phân tích
Mercer cho lớp \tr{kernel} này dựa trên các đa thức Gegenbauer. Thứ hai, bài
báo chứng minh rằng lớp \tr{kernel} được đề xuất là khá rộng, bao gồm toàn bộ
\tr{kernel} tích trong, \tr{kernel} Gaussian và một số \tr{Kernel} Tiếp tuyến
Nơ-ron. Thứ ba, tác giả đề xuất Đặc trưng Gegenbauer Ngẫu nhiên
(Random Gegenbauer Features) để xấp xỉ \tr{kernel}, đồng thời chứng minh các
bảo đảm lý thuyết như xấp xỉ phổ (spectral approximation) và bảo toàn chi phí
chiếu (projection-cost preserving). Các kết quả này cho phép áp dụng đặc trưng
ngẫu nhiên được đề xuất vào các bài toán học máy như \tr{kernel} hồi quy ridge
(ridge regression), \tr{kernel} k-means.

Ngoài ra, bài báo cũng áp dụng các kết quả lý thuyết cho \tr{kernel} tích
trong và \tr{kernel} Gaussian, qua đó cho thấy phương pháp đề xuất có thể cải
thiện so với một số phương pháp xấp xỉ \tr{kernel} trước đó. Các kết quả thực
nghiệm cũng cho thấy Đặc trưng Gegenbauer Ngẫu nhiên đạt hiệu quả tốt trong
việc xấp xỉ \tr{kernel} Gaussian.

Mối liên hệ với lý thuyết trong chương sách là rất rõ ở phần xấp xỉ
\tr{kernel} bằng ánh xạ đặc trưng. Nếu với Đặc trưng Fourier Ngẫu nhiên, ta sử
dụng định lý Bochner để viết \tr{kernel} bất biến tịnh tiến dưới dạng kỳ vọng:
\begin{equation}
  K(x,x') = \expect_{\omega \sim p}\Bigl[ \bigl[ \cos(\omega \cdot x),
    \sin(\omega \cdot x) \bigr ]^\intercal \bigl[ \cos(\omega
    \cdot x'), \sin(\omega \cdot x') \bigr ] \Bigr ],
\end{equation}
thì bài báo này phát triển một cơ chế tương tự nhưng cho lớp \tr{kernel} tổng
quát hơn. Nói cách khác, bài báo mở rộng tinh thần của Đặc trưng Ngẫu nhiên:
thay vì làm việc trực tiếp với ma trận \tr{kernel} lớn, ta xây dựng một ánh xạ
đặc trưng ngẫu nhiên sao cho tích vô hướng trong không gian đặc trưng xấp xỉ
giá trị \tr{kernel} ban đầu.

\subsubsection{Deep Learning with Kernels through
RKHM and the Perron--Frobenius Operator}

Bài báo \emph{Deep Learning with Kernels through RKHM
and the Perron--Frobenius Operator} được công bố tại NeurIPS 2023. Bài toán
chính của bài báo là xây dựng một khung lý thuyết kết hợp giữa phương pháp
\tr{kernel} và deep learning thông qua Mô-đun $C^*$ với \tr{Kernel} Tái tạo
(Hilbert Reproducing Kernel Hilbert $C^*$-Module), viết tắt là RKHM.

Trong phương pháp \tr{kernel} cổ điển, ta thường làm việc với RKHS, tức là
không gian Hilbert sinh bởi một \tr{kernel} xác định dương. Tuy nhiên, RKHS
chủ yếu phù hợp với các hàm nhận giá trị vô hướng hoặc các mở rộng
\tr{kernel} nhận giá trị \tr{vector} nhất định. Bài báo đề xuất sử dụng RKHM,
một mở rộng của RKHS dựa trên đại số $C^*$ ($C^*$-algebra), để mô hình hóa các
cấu trúc phức tạp hơn, đặc biệt là các phép hợp trong mô hình sâu.

Đóng góp chính của bài báo là đề xuất khung RKHM sâu (deep RKHM). Trong khung
này, các tầng của mô hình được mô tả thông qua các ánh xạ trong RKHM, còn sự
hợp giữa các tầng được phân tích bằng toán tử Perron--Frobenius. Bài báo cũng
đưa ra một chặn tổng quát hóa dựa trên độ phức tạp Rademacher trong bối cảnh
RKHM. Một điểm đáng chú ý là nhờ cấu trúc đại só $C^*$, chặn này có sự phụ
thuộc nhẹ hơn vào chiều đầu ra so với một số chặn trước đó.

Mối liên hệ với lý thuyết \tr{kernel} nằm ở việc bài báo mở rộng khái niệm
RKHS. Trong chương lý thuyết cơ bản, \tr{kernel} xác định một không gian RKHS
và mỗi hàm học được có thể biểu diễn thông qua \tr{kernel}. Bài báo này giữ
lại tinh thần đó nhưng thay RKHS bằng RKHM để có thể mô tả các cấu trúc hàm
phức tạp hơn. Ngoài ra, việc dùng toán tử Perron--Frobenius để phân tích phép
hợp cho thấy phương pháp \tr{kernel} không chỉ dừng lại ở mô hình một tầng, mà
còn có thể được mở rộng để nghiên cứu các mô hình sâu.

\subsubsection{Nhận xét tổng quan về xu hướng phát triển}

Từ ba bài báo trên, có thể thấy phương pháp \tr{kernel} sau năm 2021 đang phát
triển theo ba xu hướng chính.

Thứ nhất là xu hướng mở rộng quy mô. Các nghiên cứu như
\emph{Toward Large Kernel Models} cho thấy phương pháp \tr{kernel} không chỉ
là công cụ lý thuyết hoặc phù hợp với dữ liệu nhỏ, mà còn có thể được thiết kế
lại để xử lý dữ liệu lớn và tận dụng phần cứng hiện đại.

Thứ hai là xu hướng xấp xỉ \tr{kernel} hiệu quả. Các phương pháp như Đặc trưng
Gegenbauer Ngẫu nhiên tiếp tục phát triển ý tưởng Đặc trưng Ngẫu nhiên, giúp
thay thế ma trận \tr{kernel} lớn bằng các biểu diễn đặc trưng hữu hạn chiều.
Điều này làm giảm chi phí tính toán và mở rộng phạm vi áp dụng của phương pháp
\tr{kernel}.

Thứ ba là xu hướng kết nối phương pháp \tr{kernel} với học sâu. Các công trình
về \tr{Kernel} tiếp tuyến nơ-ron, học \tr{kernel} sâu (deep kernel learning)
và RKHM sâu cho thấy phương pháp \tr{kernel} đang trở thành một công cụ quan
trọng để giải thích và phân tích mạng nơ-ron sâu. Thay vì xem phương pháp
\tr{kernel} và học sâu là hai hướng tách biệt, các nghiên cứu gần đây đang cố
gắng kết hợp ưu điểm của cả hai: tính chặt chẽ lý thuyết của \tr{kernel} và
khả năng biểu diễn linh hoạt của mô hình sâu.

Nhìn chung, các hướng nghiên cứu gần đây cho thấy phương pháp \tr{kernel} vẫn
là một lĩnh vực rất quan trọng trong học máy. Trọng tâm của lĩnh vực này đang
chuyển từ các mô hình \tr{kernel} cổ điển sang các phương pháp \tr{kernel} có
khả năng mở rộng, có bảo đảm lý thuyết và có mối liên hệ sâu với học sâu hiện
đại.
